% The loop as it physically exists on this bench: six solid links and one dashed.
\begin{tikzpicture}[font=\sffamily\scriptsize, x=1cm, y=1cm]

\node[chlabel, anchor=west] at (-6.8,3.2) {the loop as it exists on this bench, drawn once so nobody has to be told twice};

% ================================================================ the host and the bus
\node[board, minimum width=28mm, minimum height=14mm] (pi) at (-5.4,1.8)
  {Raspberry Pi 4\\the outer loop};
\node[hw, text width=18mm, minimum height=7mm] (tx) at (-2.3,1.8) {transceiver};
\draw[busline] (pi.east) -- (tx.west);
\node[buslabel, anchor=south] at (-3.6,2.2) {the bus};

% ================================================================ on the node
\node[layerlabel] at (-0.6,2.9) {everything to the right of here is on the node};
\node[ctrlblk, text width=24mm, minimum height=7mm] (ctl) at (0.6,1.8) {the controller};
\node[ctrlblk, text width=24mm, minimum height=7mm, fill=fwcol] (out) at (3.4,1.8)
  {the output stage};
\node[hw, text width=20mm, minimum height=7mm] (pins) at (6.0,1.8) {two real pins};
\draw[cable] (tx.east) -- (ctl.west);
\draw[ctrlline] (ctl.east) -- (out.west);
\draw[ctrlline] (out.east) -- (pins.west);

\node[modelled, text width=24mm, minimum height=9mm] (plant) at (6.0,0.2)
  {the plant model\\\textbf{not on this bench}};
\draw[ctrlline] (pins.south) -- (plant.north);

\node[hw, text width=24mm, minimum height=7mm] (qgen) at (3.4,-1.3) {quadrature out};
\node[sensorblk, text width=24mm, minimum height=7mm] (enc) at (0.6,-1.3)
  {the encoder input};
\draw[ctrlline] (plant.south) -- (6.0,-1.3) -- (qgen.east);
\draw[cable] (qgen.west) -- (enc.east);
\draw[fbline] (enc.north) -- (ctl.south);

\node[buslabel, anchor=north] at (2.0,-1.75) {two pins to two pins};
\node[chlabel, anchor=west] at (1.0,0.3) {the measurement, once per period};

% ================================================================ the count
\node[regcell, fill=usrcol, minimum width=54mm, text width=52mm, inner sep=3pt,
      anchor=north west, align=left] at (-6.8,-2.4)
  {\textbf{Six solid links:} bus in, controller, output stage, real pins, real
   quadrature pulses, real encoder input.\\[2pt]
   \textbf{One dashed:} the mechanics, which is software.};

\node[regcell, fill=extcol, minimum width=54mm, text width=52mm, inner sep=3pt,
      anchor=north west, align=left] at (0.2,-2.4)
  {The feedback path contains real edges, real capture hardware and real
   interrupt latency. What it does not contain is friction, backlash, compliance
   or a load, and chapter 7 said which of those the model has.};

\node[umlnote, text width=114mm, anchor=north west] at (-6.8,-4.6)
  {This arrangement is more honest than a pure simulation and less capable than
   a motor, and both halves of that sentence belong in the portfolio text. The
   timing, the peripherals and the arithmetic are exercised exactly as they would
   be with a motor attached. The mechanics are a model whose parameters somebody
   chose, and a controller tuned against it is tuned against that choice.};
\end{tikzpicture}
